
\begin{proposition}[Every reduced all-Hopf support has one cycle, with an exact algebraic-dimension-zero classification]\label{cdp:st:C.2.UN.54}
Let $f:U\to\Delta$ be a proper degeneration of connected complex
two-tori with smooth total space, whose full central fibre
$S=\bigcup_{v\in V}H_v$ is reduced and simple normal crossing.
Suppose that every $H_v$ is a primary or secondary Hopf surface.
Then every connected double curve is elliptic, there are no triple
points, and the connected incidence graph has exactly one independent
cycle.

For an oriented edge $e$ from $v$ to $w$, let $C_e$ be its actual
double curve and let
\[
 r_{v,e}:H^1(H_v,\OO_{H_v})\longrightarrow H^1(C_e,\OO_{C_e})
                                                        \tag{C.HR.1}\label{cdp:eq:C.HR.1}
\]
be restriction through the specified inclusion (and through the
specified gluing map at the other endpoint). Every such map is a
nonzero map between one-dimensional complex vector spaces. The
weighted incidence map
\[
 d_1:\bigoplus_vH^1(H_v,\OO_{H_v})
       \longrightarrow\bigoplus_eH^1(C_e,\OO_{C_e}),\qquad
 (d_1a)_e=r_{v,e}a_v-r_{w,e}a_w
                                                        \tag{C.HR.2}\label{cdp:eq:C.HR.2}
\]
has rank $|V|-1$. In any choices of nonzero bases, the product of
$r_{v,e}/r_{w,e}$ around the oriented cycle is one. This product
retains the actual restriction maps; it is unchanged by changing
the chosen bases. The full coherent dimensions are
\[
                    h^q(S,\OO_S)=(1,2,1),\qquad q=0,1,2.
                                                        \tag{C.HR.3}\label{cdp:eq:C.HR.3}
\]

The one-dimensional space $H^0(S,\omega_S)$ has a more precise
description. Its nonzero section is represented by logarithmic
two-forms on the components. Their residues are nonzero on exactly
the cycle edges and zero on every edge of an attached tree. The
forms themselves are nonzero on exactly the cycle components.
For a canonical evaluation generator with vertical orders $k_v$,
the minimum of the $k_v$ is therefore attained exactly at the
cycle components.

In particular, if every $H_v$ has algebraic dimension zero, the
graph has no attached trees and every component is a diagonal Hopf
surface. There are at least two components. This gives the following
complete criterion for this entire reduced branch, without initially
assuming either diagonality or a cycle. A specified compact reduced
SNC surface all of whose components are algebraic-dimension-zero
Hopf surfaces occurs as the full fibre of a smooth-total torus
degeneration precisely when its components are diagonal, its graph
is a cycle, its actual axis normal lines are dual, and the integral
edge product $M$ of (C.HC.3)--(C.HC.6) is the identity. Every such
specified support has the compact realization of \ref{cdp:st:C.2.UN.48}, with
all support Euler characteristics zero.
\end{proposition}

\begin{proof}
\par\medskip\noindent\textbf{The curves and the actual restriction maps.}
A Hopf surface has $H^2(H,\Q)=0$. For an irreducible compact curve
$C\subset H$, adjunction gives $p_a(C)=1$. A map from $\PP^1$ to
$H$ lifts to its universal cover $\C^2\setminus\{0\}$, and its two
coordinate functions are constant. The normalization of $C$ hence
cannot be rational. The equality
$p_a(C)=g(C^\nu)+\delta(C)$ now gives $g(C^\nu)=1$ and
$\delta(C)=0$. Thus $C$ is smooth elliptic. Two distinct such
curves have intersection number zero, since $H^2(H,\Q)=0$;
positivity of every local intersection forces them to be disjoint.
These are the same explicit curve comparisons used in \ref{cdp:st:C.2.UN.37}.
It follows at once that the present support has only elliptic
double curves and no triple point.

We prove nonvanishing in \eqref{cdp:eq:C.HR.1}, including secondary surfaces.
First let $H'=(\C^2\setminus\{0\})/\langle F\rangle$ be a primary
Hopf surface. The constant additive deck cocycle $1$ represents a
nonzero class in $H^1(H',\OO_{H'})$. Indeed, a trivialization would
be a holomorphic function $b$ on $\C^2\setminus\{0\}$ satisfying
$b(Fz)-b(z)=1$. Hartogs extension to the origin and evaluation there
give a contradiction. Since $h^1(H',\OO_{H'})=1$, this class spans
the group.

For every elliptic curve $C'\subset H'$, its lift can have finitely
many branches permuted by $F$. Choose a positive power $F^d$
preserving one branch and removing any root-of-unity discrepancy
in the diagonal multipliers. The actual curve list used in
\ref{cdp:st:C.2.UN.38} and \ref{cdp:st:C.2.UN.53} gives its normalized parameter $\C^*$,
with $F^d$ acting by $u\mapsto\lambda u$, where
$0<|\lambda|<1$. For an axis this is its displayed diagonal
multiplier; for a resonant binomial curve $z^q=cw^p$ it follows
from $(z,w)=(a u^p,b u^q)$. The scalar and nonlinear-axis cases
are included. The actual map
$\C^*/\langle F^d\rangle\to C'$ induced by this inclusion is a
finite unramified elliptic covering. Pulling the constant deck
cocycle $1$ to the subgroup generated by $F^d$ gives the cocycle
$d$. Thus, if the original restriction were zero, its pullback
would give a holomorphic function on $\C^*$ satisfying
\[
                         b(\lambda u)-b(u)=d.          \tag{C.HR.4}\label{cdp:eq:C.HR.4}
\]
In its Laurent expansion on $\C^*$ the constant coefficient on
the left is zero, whereas $d>0$. This proves nonvanishing without
assuming that $F$ fixes every lifted branch individually.

For a secondary $H$, take the actual finite unramified primary
cover $\pi:H'\to H$ of \ref{cdp:st:C.2.UN.43}. A component $C'$ of the inverse
image of $C$ maps by a finite unramified elliptic covering to $C$.
The trace identities
$\operatorname{Tr}\circ\pi^*=(\deg\pi)\operatorname{id}$, on $H$
and on this curve covering, show that both pullbacks on $H^1(\OO)$
are injective. The surface pullback is an isomorphism of
one-dimensional spaces. Naturality of the restriction square,
together with the primary calculation, proves that the original
$H^1(H,\OO_H)\to H^1(C,\OO_C)$ is nonzero. Thus every coefficient
in \eqref{cdp:eq:C.HR.2} is nonzero, with the actual gluing maps included.

\par\medskip\noindent\textbf{Normalization, the graph, and its exact cycle product.}
The no-triple SNC equation gives the exact normalization sequence
\[
 0\longrightarrow\OO_S
 \longrightarrow\bigoplus_v\OO_{H_v}
 \xrightarrow{\ d\ }\bigoplus_e\OO_{C_e}
 \longrightarrow0,                                    \tag{C.HR.5}\label{cdp:eq:C.HR.5}
\]
where $d$ is the difference of the actual two restrictions. On
$H^0$ it is the ordinary signed incidence matrix $d_0$. The
connectedness of $S$ gives $\operatorname{rank}d_0=|V|-1$.
Since $H^2(H_v,\OO_{H_v})=0$, its long exact sequence gives
\[
 \begin{gathered}
 0\longrightarrow\operatorname{coker}d_0
 \longrightarrow H^1(S,\OO_S)
 \longrightarrow\ker d_1\longrightarrow0,\\
             H^2(S,\OO_S)\simeq\operatorname{coker}d_1.
 \end{gathered}                                        \tag{C.HR.6}\label{cdp:eq:C.HR.6}
\]
Here $\dim\operatorname{coker}d_0$ is the first Betti number of
the incidence graph.

Choose a spanning tree and a nonzero value at one vertex.
Its edge equations $r_{v,e}a_v=r_{w,e}a_w$ successively determine
nonzero values at all other vertices. Thus the tree rows of $d_1$
are independent and have a one-dimensional kernel. In particular
if the whole graph were a tree, $d_1$ would be surjective and
$H^2(S,\OO_S)=0$.

The map $f$ is flat: multiplication by its nonzero local parameter
is injective in each regular local ring of the smooth total space,
and torsion-free modules over the base discrete valuation ring
are flat. Upper semicontinuity for this proper flat family gives
$h^2(S,\OO_S)\ge h^2(F,\OO_F)=1$. The central fibre is the
reduced $S$ by hypothesis, so no thickened-fibre cohomology has
been substituted here. The tree case is impossible.

\ref{cdp:st:C.2.UN.44} gives a surjection from $\pi_1(S)$ to the incidence
graph group, and $\pi_1(S)$ is a quotient of $\Z^4$ because the
full fibre is reduced. The free graph group is therefore abelian
and has rank at most one. Together with the preceding paragraph
this proves that the graph is unicyclic and $|E|=|V|$.

There is now just one row left after the spanning-tree equations.
Writing its endpoints using the successive nonzero values shows
that it vanishes precisely when the product of the actual ratios
$r_{v,e}/r_{w,e}$ around the cycle is one. If it did not vanish,
$d_1$ would have rank $|V|$ and again $H^2(S,\OO_S)=0$.
The product is consequently one, the rank is $|V|-1$, and
\eqref{cdp:eq:C.HR.6} gives exactly \eqref{cdp:eq:C.HR.3}. A vertex-basis change cancels
between the incoming and outgoing ratios; an edge-basis change
cancels inside its ratio. This proves the asserted independence
without discarding the analytic restriction coefficients.

\par\medskip\noindent\textbf{The dualizing form and the canonical-order minimum.}
Let $D_v$ be the disjoint union of the double curves on $H_v$.
The residue exact sequence is
\[
 0\longrightarrow K_{H_v}
 \longrightarrow K_{H_v}(D_v)
 \xrightarrow{\operatorname{Res}}\bigoplus_{e\ni v}K_{C_e}
 \longrightarrow0.                                    \tag{C.HR.7}\label{cdp:eq:C.HR.7}
\]
Since $H^0(K_{H_v})=0$, a logarithmic two-form is determined by
its residues. Serre duality identifies the connecting map from
these residues to $H^1(K_{H_v})$ with the dual of the actual
restriction maps \eqref{cdp:eq:C.HR.1}. The signs may be fixed explicitly:
in a double chart $xy=0$ the two residues of a dualizing form
sum to zero. Orienting the edge takes its residue on the tail
and the negative on the head. Thus the global residue vector of
a section of $\omega_S$ belongs exactly to $\ker d_1^\vee$.
Conversely a vector in that kernel lifts uniquely by \eqref{cdp:eq:C.HR.7} on
each component, and the opposite-residue rule glues these forms
to a section of $\omega_S$. These are inverse maps.

At a leaf of an attached tree, the corresponding row of
$d_1^\vee$ contains only its one nonzero coefficient. The
residue on that edge is zero. Delete that leaf and repeat.
Every tree-edge residue is zero. On the remaining cycle each
vertex relates its two adjacent residues by nonzero coefficients;
its one-dimensional kernel has a nonzero residue on every cycle
edge. On an off-cycle component all residues vanish, so the
form is an ordinary holomorphic two-form and is zero. On a cycle
component its two nonzero residues make it nonzero. Vanishing
residue along any other boundary curve removes the logarithmic
pole there, directly from the local expression
$a(y,z)\,dy\wedge dz/y$.

Finally take the canonical evaluation section of \ref{cdp:st:C.2.UN.38} and
write its vertical divisor as $\sum k_vH_v$. Put $k_0=\min k_v$.
Because $S$ is the full reduced divisor of $t$, division by
$t^{k_0}$ is an actual holomorphic section of the relative
canonical line. Its restriction to $S$ is a nonzero section of
$\omega_S$. At an edge joining $v,w$, the explicit adjunction
calculation, with $t=xy$, is
\[
 \frac{x^{k_v-k_0}y^{k_w-k_0}\,dx\wedge dy\wedge dz}{xy}
 \quad\longmapsto\quad
 y^{k_w-k_0-1}\,dy\wedge dz\quad\text{on }x=0
                                                        \tag{C.HR.8}\label{cdp:eq:C.HR.8}
\]
when $k_v=k_0$, times a holomorphic unit; when $k_v>k_0$ the
restriction on that component is zero. The local logarithmic
expression is the usual restriction of the relative dualizing
line, or equivalently the residue of the displayed meromorphic
three-form. It follows that its nonzero components are exactly
those with $k_v=k_0$. The preceding residue calculation identifies
these with the cycle components, as claimed.

\par\medskip\noindent\textbf{The full algebraic-dimension-zero branch.}
The actual primary and secondary curve computations in
\ref{cdp:st:C.2.UN.38} and \ref{cdp:st:C.2.UN.43} give two possibilities in algebraic
dimension zero: a nonlinear Hopf surface has just one elliptic
curve, and a diagonal Hopf surface has just its two distinct
axis elliptic curves. The diagonal finite action cannot exchange
them: such an exchange would force equal positive powers of the
two contracting multipliers, giving a nonconstant meromorphic
invariant and algebraic dimension one. Thus every vertex of the
present graph has degree at most two.

A connected unicyclic finite graph with all degrees at most two
is a cycle and has every degree equal to two. Indeed its equality
$|E|=|V|$ gives average degree two, so no smaller degree is possible.
Consequently there is no nonlinear component and no attached
tree. There is no loop at a single vertex, since the components
of the specified SNC surface are smooth embedded components:
such a loop would be a self-identification, not a double curve
between distinct components. Hence there are at least two
components.

For this now established diagonal cycle, the actual normal-product
trivialization of $t=xy$ gives the dual-normal maps
(C.HC.3). Removing the zero sections gives the actual open groups
$G_v=\C^2/\Lambda_v$ and the affine logarithmic edge maps
(C.HC.5). Their integral lattice product is precisely $M$, and
\ref{cdp:st:C.2.UN.46} proves that an actual torus smoothing forces $M=I$.
Conversely the specified axis maps and any compatible dual-normal
maps with $M=I$ give the proper plumbing of \ref{cdp:st:C.2.UN.46}; the actual
varying-period comparison in \ref{cdp:st:C.2.UN.48} extends it to a compact
smooth threefold over a compact curve. That theorem retains every
component and every specified gluing constant, so the central
support is isomorphic to the specified $S$, rather than merely
having the same graph. Its other supports have Euler zero.

The exact matrix calculation in
\path{code/E_hopf_graph_reduction.py} checks the spanning-tree
elimination, the cycle product, both coherent dimensions, and the
support of the dual residue vector on several branched graphs.
The argument above proves each comparison for every finite graph.
\end{proof}


